\documentclass[11pt, a4paper, oneside]{ctexart}
\usepackage{indentfirst, amsmath, amsthm, amssymb, graphicx, float, subfigure, geometry, setspace, xfrac, tikz-cd, enumerate}
\usepackage[bookmarks=true, colorlinks, citecolor=blue, linkcolor=black]{hyperref}
\ctexset{
    section = {
        format = {\Large\bfseries},
    }
}

% 导言区

\title{Note 03}
\author{zero-point\_}
\setlength{\parindent}{4ex}
\pagestyle{plain}
\geometry{left=2.54cm, right=2.54cm, top=3.18cm, bottom=3.18cm}
\setstretch{1.5}

\newtheorem{theorem}{定理}[section]
\newtheorem{corollary}[theorem]{推论}
\newtheorem{definition}[theorem]{定义}
\newtheorem{example}[theorem]{例}
\newtheorem{lemma}[theorem]{引理}
\newtheorem{proposition}[theorem]{命题}
\newtheorem{remark}[theorem]{注}

\newcommand{\C}{\mathbb{C}}
\newcommand{\N}{\mathbb{N}}
\newcommand{\Q}{\mathbb{Q}}
\newcommand{\R}{\mathbb{R}}
\newcommand{\T}{\mathbb{T}}
\newcommand{\Z}{\mathbb{Z}}
\newcommand{\cA}{\mathcal{A}}
\newcommand{\cB}{\mathcal{B}}
\newcommand{\cF}{\mathcal{F}}
\newcommand{\cO}{\mathcal{O}}
\newcommand{\Id}{\mathrm{Id}}
\newcommand{\re}{\mathrm{Re}}
\renewcommand{\top}{\mathrm{top}}
\newcommand{\tr}{\mathrm{tr}}
\newcommand{\abs}[1]{\left|{#1}\right|}
\newcommand{\norm}[1]{\left\Vert{#1}\right\Vert}
\renewcommand{\v}[1]{\mathbf{#1}}

\begin{document}
	\maketitle
    \section{内容安排}
    \begin{itemize}
        \item 结构稳定性简介
        \item 周期点拓扑不变量
        \item 拓扑熵定义
        \item 拓扑熵计算
    \end{itemize}

    \section{结构稳定性}
    结构稳定性与讨论班的主线关系不大，但可以直观感受到动力系统的性质，因此我们简要介绍一下.
    \begin{definition}[$C^m$ 结构稳定性]
        一个 $C^r$ 映射 $f$ 称为 $C^m$ 结构稳定的（$1\leq m\leq r$），如果存在 $C^m$ 拓扑下的 $f$ 的一个邻域 $U$，使得任意 $g\in U$ 都与 $f$ 拓扑共轭.
    \end{definition}
    \begin{definition}[$C^m$ 强结构稳定性]
        一个 $C^r$ 映射 $f$ 称为 $C^m$ 强结构稳定的（$1\leq m\leq r$），如果存在 $C^m$ 拓扑下的 $f$ 的一个邻域 $U$，使得任意 $g\in U$ 都与 $f$ 拓扑共轭（结构稳定），并且共轭映射 $h=h_g$ 在 $g\stackrel{C^m}{\to}f$ 时满足 $h_g,h_g^{-1}$ 一致趋于 $\Id$.
    \end{definition}
    \begin{definition}[拓扑稳定性]
        一个 $C^r$ 微分同胚 $f$ 称为拓扑稳定的，如果存在 $C^0$ 拓扑（一致拓扑）下的 $f$ 的一个邻域 $U$，使得对任意 $g\in U$，若 $g$ 是同胚，则与 $f$ 拓扑半共轭（$f$ 是 $g$ 的因子）.
    \end{definition}
    \begin{remark}
        \begin{itemize}
            \item $C^0$ 扰动比较多样，我们不好去强求拓扑共轭，因此退而求其次要求半共轭.
            \item 若是局部微分同胚，同样的定义也可以良定义.
        \end{itemize}
    \end{remark}

    上述稳定性是对离散迭代而言的，对于连续流，拓扑共轭性质太好，因此我们退而求其次，使用轨道等价性：
    \begin{definition}[轨道等价性]
        $\varphi^t,\psi^t$ 分别是 $M,N$ 上的 $C^r$ 流，称二者 $C^m$ 轨道等价，若存在一个 $C^m$ 微分同胚 $h:M\to N$，使得 $h(\varphi^t(x))=\psi^{\tau(t,x)}(h(x))$，其中 $\tau:\R\times M\to\R$ 是一个连续函数，并且对于任意 $x\in M$，$\tau(\cdot,x)$ 是一个单调函数（$\psi^{\tau(t,x)}(x)$ 也称为 $\psi$ 的时间变化）.
    \end{definition}
    对应的，我们也给出轨道因子的定义：
    \begin{definition}[轨道因子]
        $\varphi^t,\psi^t$ 分别是 $M,N$ 上的 $C^r$ 流，称 $\psi^t$ 是 $\varphi^t$ 的 $C^m$ 轨道因子，若存在一个 $C^m$ 满射同胚 $h:M\to N$，使得 $\varphi^t$ 的轨道被映成 $\psi^t$ 的轨道.
    \end{definition}

    \begin{definition}[流的 $C^m$（强）结构稳定性]
        一个 $C^r$ 流 $\phi^t$ 称为 $C^m$ 结构稳定的（$1\leq m\leq r$）（对应称为强结构稳定的），如果存在 $C^m$ 拓扑下的 $\phi^t$ 的一个邻域 $U$，使得任意 $\psi^t\in U$ 都与 $\phi^t$ 轨道等价.
    \end{definition}
    \begin{definition}[流的拓扑稳定性]
        一个 $C^r$ 流 $\varphi^t$ 称为拓扑稳定的，如果存在 $C^0$ 拓扑下的 $\varphi^t$ 的一个邻域 $U$，使得 $\varphi^t$ 是任意 $\psi^t\in U$ 的轨道因子.
    \end{definition}

    \section{周期点拓扑不变量}
    类似于同调理论等，我们给出一个等价关系，就希望可以找到等价关系下的一些不变量，并希望这些不变量可以尽可能区分更多不等价的元素. 既然是动力系统课程，我们显然会关心函数在长期迭代下是否可以定义出一个量，且这个量在拓扑共轭下不变. 以下我们介绍周期点的一个拓扑不变量作为引入.
    \begin{definition}[周期点渐近增长指数]
        定义周期点渐近增长指数 $p(f)=\varlimsup\limits_{n\to\infty}\frac{\log\max(P_n(f),1)}{n}$.
    \end{definition}
    \begin{remark}
        \begin{itemize}
            \item 由于对特定 $n$ 可能没有周期点，因此我们取 $P_n(f)$ 和 $1$ 之间的较大值.
            \item 根据此前几个样例的分析，周期点数经常是指数增长的，因此我们取的是指数的计算值. 但是这不总正确，例如 $p(f)=0$ 时，我们也可以考虑多项式次数 $\varlimsup\limits_{n\to\infty}\frac{\log\max(P_n(f),1)}{\log n}$，等等.
            \item $p(f)<\infty$ 时，我们实际上可以把每个 $n$ 的周期点数信息集中在一个 $\zeta$ 函数 $\zeta_f(z)=\exp(\sum\limits_{n=1}^{\infty}\frac{P_n(f)}{n}z^n)$，这一级数在 $\abs{z}<e^{-p(f)}$ 时收敛（请证明），边界上必有奇点. 很多时候，奇点只是极点，因此可能可以解析延拓到 $\C$ 上的亚纯函数（请计算之前介绍的几个例子来验证），因此 $\zeta_f$ 的各个指标（如极点、零点、留数等）都可以提供额外的拓扑不变量.
        \end{itemize}
    \end{remark}

    对于流，由于是连续情形，我们很难用“等于某个周期”的条件进行计数，而若周期轨道数指数级增长，我们可以期望使用“小于等于某个周期”的条件进行计数，因此我们定义如下的渐近增长指数：
    \begin{definition}[周期轨道渐近增长指数]
        定义周期轨道渐近增长指数 $p(\varphi^t)=\varlimsup\limits_{T\to\infty}\frac{\log\max(P_T(\varphi^t),1)}{T}$，其中 $P_T(\varphi^t)$ 指的是周期不大于 $T$ 的周期轨道数.
    \end{definition}
    \begin{remark}
        \begin{itemize}
            \item 类似的，我们可以定义 $\zeta$ 函数 $\zeta_{\varphi^t}(z)=\prod\limits_{\gamma}{(1-e^{-zl(\gamma)})}^{-1}$，这里 $l(\gamma)$ 是 $\gamma$ 的最小正周期，且要求 $\gamma$ 不是不动点轨道. 收敛性上，若 $p(\varphi^t)$ 可以良定义，则 $\re(z)>p(\varphi^t)$ 时无穷乘积收敛（请验证），并且在边界上有奇点.
            \item 这两个定义是与周期强相关的，因此如果两个流拓扑共轭，它们就是不变量；但如果两个流只是轨道等价，它们就不一定是不变量了. 然而，在特殊情况下我们有以下定理：
        \end{itemize}
    \end{remark}

    \begin{theorem}
        设 $X,Y$ 是紧空间，$\varphi^t:X\to X$ 和 $\psi^t:Y\to Y$ 是无不动点的连续流，且二者轨道等价. 若 $p(\psi^t)=0$，则 $p(\varphi^t)=0$.
    \end{theorem}
    \begin{proof}
        设 $h:X\to Y$ 是轨道等价中的同胚，$\tau(t,x)$ 是轨道等价中的时间变化函数，取切片 $\alpha(x)=\tau(1,x)$，则由于 $X$ 的紧性知 $\alpha$ 有界 $M$，则 $\varphi^t$ 长度（这里指的是运行时间）为 $1$ 的周期轨道在 $\psi^t$ 中的长度不大于 $M$，因此当 $T$ 足够大时，$P_T(\varphi^t)\leq P_{(T+1)M}(\psi^t)\leq P_{2MT}(\psi^t)$（注意向上取整的事情），从而 $p(\psi^t)\geq \frac{1}{2M}p(\varphi^t)$，特别的，$p(\psi^t)=0$ 时 $p(\varphi^t)=0$.
    \end{proof}

    \section{拓扑熵的定义}
    接下来我们引入最重要的不变量定义：拓扑熵. 虽然它不依赖度量，但使用度量语言可以直观给出它的 Bowen-Dinaburg 定义.

    首先，我们给出基于 $f$ 的一种度量定义，并探讨度量空间的几个重要概念.

    \begin{definition}[映射诱导度量]
        设 $f:X\to X$ 是紧度量空间 $X$ 上的连续函数，$X$ 度量为 $d$，则定义 $f$ 在 $n$ 次迭代中诱导的度量 $d_n^f(x,y)=\max\limits_{0\leq i\leq n-1}d(f^i(x),f^i(y))$. 我们记开球 $B_f(x,\varepsilon,n)=\{y\in X\mid d_n^f(x,y)<\varepsilon\}$.
    \end{definition}
    \begin{remark}
        显然，固定两点时，$d_n^f$ 单调递增；
    \end{remark}
    \begin{definition}[覆盖与分离]
        给定集合 $E\subset X$. 若 $X\subset \bigcup\limits_{x\in E}B_f(x,\varepsilon,n)$，则称 $E$ 是 $X$ 的一个 $(n,\varepsilon)$-覆盖. 若 $E$ 中任意两点 $x,y$ 满足 $d_n^f(x,y)\geq\varepsilon$，则称 $E$ 是 $X$ 的一个 $(n,\varepsilon)$-分离集.

        对应的，我们定义 $S_d(f,\varepsilon,n)$ 是 $(n,\varepsilon)$-覆盖的最小元素个数，$N_d(f,\varepsilon,n)$ 是 $(n,\varepsilon)$-分离集的最大元素个数；我们另外定义 $D_d(f,\varepsilon,n)$ 是 $X$ 满足以下条件的覆盖集 $\cF$ 元素个数的最小值：$\cF$ 中每个元素 $A$ 的半径都不大于 $\varepsilon$.
    \end{definition}
    \begin{remark}
        \begin{itemize}
            \item $S_d$ 与 $D_d$ 有很好的相容性：$S_d(f,\varepsilon,n)\leq D_d(f,\varepsilon,n)\leq S_d(f,\frac{\varepsilon}{2},n)$（请证明）. 但是二者存在不同：$D_d$ 允许非开集的存在，因此在定义上性质更好，而 $S_d$ 只能用开球，因此在计算上性质更好，我们接下来就可以看到二者的不同.
            \item $N_d$ 与 $S_d$ 也有以下不等式：$S_d(f,\varepsilon,n)\leq N_d(f,\varepsilon,n)\leq S_d(f,\frac{\varepsilon}{2},n)$（请证明）.
        \end{itemize}
    \end{remark}

    我们先给出使用 $S_d$ 的拓扑熵定义，再考虑与其他定义的相容性.
    \begin{definition}[拓扑熵]
        定义 $h_d(f,\varepsilon)=\varlimsup\limits_{n\to\infty}\frac{1}{n}\log S_d(f,\varepsilon,n)$，则由于 $h_d$ 对于 $\varepsilon\to 0$ 过程的单调不减性，以下极限良定义：$h(f)=h_{\top}(f)=h_d(f)=\lim\limits_{\varepsilon\to 0}h_d(f,\varepsilon)$.
    \end{definition}
    我们立刻给出度量依赖性的说明：
    \begin{proposition}[拓扑熵不依赖度量]
        设度量 $d,d'$ 给出同一个拓扑，则 $h_d(f)=h_{d'}(f)$.
    \end{proposition}
    \begin{proof}
        设 $D_{\varepsilon}=\{(x_1,x_2)\mid d(x_1,x_2)\geq \varepsilon\}\subset X\times X$，这是 $X\times X$ 上的紧集（因为是闭的）. 由于 $d'$ 在 $X\times X$ 上连续，则在紧集 $D_{\varepsilon}$ 上可以取得最小值 $\delta(\varepsilon)>0$（为什么？）于是 $d'(x_1,x_2)<\delta(\varepsilon)\Rightarrow d(x_1,x_2)<\varepsilon$，于是对 $d_n^f$ 同理，于是 $S_{d'}(f,\delta(\varepsilon),n)\geq S_d(f,\varepsilon,n)$，于是 $h_{d'}(f,\delta(\varepsilon))\geq h_d(f,\varepsilon)$，两边令 $\varepsilon\to 0$ 则知 $h_{d'}(f)\geq h_d(f)$，由对称性即证.
    \end{proof}
    \begin{corollary}[拓扑熵是不变量]
        拓扑熵是拓扑共轭的不变量.
    \end{corollary}
    \begin{proof}
        设 $f:X\to X,g:Y\to Y$ 通过 $h:X\to Y$ 拓扑共轭，则固定 $X$ 上的度量 $d$，设 $d'(y_1,y_2)=d(h^{-1}(y_1),h^{-1}(y_2))$ 是拉回度量（请验证它是度量），则 $h$ 是等距同构，即证.
    \end{proof}

    从上述 $S_d,D_d$ 与 $D_d,N_d$ 的不等式，我们也容易通过 $D_d$ 或 $N_d$ 类似地定义拓扑熵（$N_d$ 要取下极限），而且它们与 $S_d$ 定义的拓扑熵相合. 接下来我们说明 $D_d$ 的良好定义性质. 我们先不加证明地给出一个关键引理（教材命题 9.6.4）：
    \begin{lemma}[Fekete]
        对数列 $\{a_n\}$，若 $\exists\; L\in\R,\; \exists\; k\in\N,\; \forall\; m,n\in\N,\; a_{m+n}\leq a_n+a_{m+k}+L$，则 $\lim\limits_{n\to\infty}\frac{a_n}{n}\in \R\cup\{-\infty\}$ 存在.
    \end{lemma}
    \begin{lemma}
        任取 $\varepsilon>0$，$\lim\limits_{n\to\infty}\frac{1}{n}\log D_d(f,\varepsilon,n)$ 存在.
    \end{lemma}
    \begin{proof}
        我们只需证明 $D_d(f,\varepsilon,n)$ 的次乘性，即 $D_d(f,\varepsilon,n+m)\leq D_d(f,\varepsilon,n)D_d(f,\varepsilon,m)$，则 Fekete 引理即可得到结论. 要证明不等式，只需注意到若 $A,B\subset X$ 分别是 $d_n^f,d_m^f$-半径小于 $\varepsilon$ 的集合，则 $A\cap f^{-n}(B)$ 是 $d_{n+m}^f$-半径小于 $\varepsilon$ 的集合.
    \end{proof}
    \begin{remark}
        若按 $\varliminf\limits_{n\to\infty}\frac{1}{n}\log S_d(f,\varepsilon,n)$ 来定义 $h_d(f,\varepsilon)$，再定义 $h_d(f)$，则仍然与 $D_d$ 定义的 $h_d(f)$ 相合，因此其也与当前用上极限定义的 $h_d(f)$ 相合.
    \end{remark}

    拓扑熵不依赖度量，因此定义可以不使用度量，接下来我们给出 Adler-Konheim-McAndrew 的拓扑定义.
    \begin{definition}[拓扑熵]
        给定紧集上的连续映射 $f:X\to X$，设 $\alpha$ 是 $X$ 的有限开覆盖，定义 $H(\alpha)=\log N(\alpha)$，其中 $N(\alpha)$ 是 $\alpha$ 的最小子覆盖的开集数；定义 $\alpha_n=\bigvee\limits_{i=0}^{n-1}f^{-i}(\alpha)$，其中集族 $\alpha\vee\beta=\{A\cap B\mid A\in \alpha,B\in\beta\}$. 定义 $h(f,\alpha)=\lim\limits_{n\to\infty}\frac{1}{n}H(\alpha_n)$，再定义拓扑熵 $h(f)=\sup\limits_{\alpha}h(f,\alpha)$.
    \end{definition}
    \begin{remark}
        ${\{H(\alpha_n)\}}_{n\in\N}$ 有次可加性，因此定义才能成立.
    \end{remark}

    \begin{theorem}[相合性]
        在紧度量空间上，两种定义是相合的.
    \end{theorem}
    \begin{proof}
        拓扑定义 $\leq$ 度量定义：只需证明每个给定的 $\alpha$ 都存在 $\varepsilon>0$ 使得 $h(f,\alpha)\leq h(f,\varepsilon)$. 由 Lebesgue 数引理，取 $\varepsilon$ 为 $\alpha$ 的 Lebesgue 数，使得任意半径为 $\varepsilon$ 的开球都包含在 $\alpha$ 的某个元素中. 接下来设 $E$ 是一个 $(n,\varepsilon)$-覆盖，则 $\forall\; x\in E$，由定义知 $B_f(x,\varepsilon,n)=\bigcap\limits_{i=0}^{n-1}f^{-i}(B(f^i(x),\varepsilon))$. 由 Lebesgue 数知 $f^{-i}(B(f^i(x),\varepsilon))\subset A_i\in\alpha$，则 $B_f(x,\varepsilon,n)\subset A_x\in\alpha_n$，由于 $N(\alpha_n)$ 是 $\alpha_n$ 的最小子覆盖的开集数，而 $\{B_f(x,\varepsilon,n)\mid x\in E\}$ 可以覆盖 $X$，于是 $S_d(f,\varepsilon,n)\geq N(\alpha_n)$，于是两边对 $n$ 取上极限即证.
        
        度量定义 $\leq$ 拓扑定义：同样我们要证明 $\forall\; \varepsilon>0$，存在有限开覆盖 $\alpha$ 使得 $h(f,\varepsilon)\leq h(f,\alpha)$. $\alpha$ 与 $n$ 无关，所以我们先用 $\varepsilon$ 构造出 $\alpha$：取 $\alpha$ 是 $\{B(x,\frac{\varepsilon}{2})\mid x\in X\}$ 的最小开覆盖，则 $\alpha_n$ 中的任意元素 $U=\bigcap\limits_{i=0}^{n-1}f^{-i}(B(x_i,\frac{\varepsilon}{2}))$，任取 $(n,\varepsilon)$-分离集 $E$，若 $E\cap U\supset \{x,y\}$（包含不止两点），则 $d(f^i(x),f^i(y))\leq d(f^i(x),x_i)+d(f^i(y),x_i)<\varepsilon$，于是 $d_n^f(x,y)<\varepsilon$，矛盾. 从而 $E\cap U$ 至多只有一点，于是 $N(\alpha_n)\geq \abs{E}$ 对任意 $(n,\varepsilon)$-分离集 $E$ 都成立，于是 $N(\alpha_n)\geq N_d(f,\varepsilon,n)$，两边对 $n$ 取上极限即证（这里 $h(f,\varepsilon)$ 改为使用 $N_d$ 定义，虽然与 $S_d$ 定义下的结果不同，但最后计算拓扑熵的极限结果相同）.
    \end{proof}

    最后，我们说明，上述对离散迭代的定义可以通过修改度量的方式迁移到流上：度量修改为 $d_t^f(x,y)=\max\limits_{0\leq t\leq T}d(\varphi^t(x),\varphi^t(y))$，这里由于空间紧而最大值良定义.

    \section{拓扑熵的基本性质}
    接下来我们给出拓扑熵的一些基本性质.
    
    \begin{proposition}
        若 $g:Y\to Y$ 是 $f:X\to X$ 的因子，则 $h(g)\leq h(f)$.
    \end{proposition}
    \begin{proof}
        设 $h\circ f=g\circ h$，则由于一切在紧空间上定义，则 $h$ 一致连续，则 $\forall\; \varepsilon>0,\; \exists\; \delta>0$ 使得 $d_X(x_1,x_2)<\delta\Rightarrow d_Y(h(x_1),h(x_2))<\varepsilon$，则 $S_{d_X}(f,\delta,n)\geq S_{d_Y}(g,\varepsilon,n)$.
    \end{proof}

    \begin{proposition}
        \begin{enumerate}[(1)]
            \item 设 $\Lambda$ 是 $f$ 的闭不变集，则 $h(f\mid_{\Lambda})\leq h(f)$.
            \item 设 $X$ 是若干闭不变子集 $\Lambda_i,i=1,\ldots,m$ 的并，则 $h(f)=\max\limits_{1\leq i\leq m}h(f\mid_{\Lambda_i})$.
            \item $h(f^m)=\abs{m}h(f)$（当 $f$ 可逆时 $m<0$ 对应拓扑熵良定义）.
            \item $h(f\times g)=h(f)+h(g)$，这里 $f:X\to X,g:Y\to Y$ 则 $f\times g:X\times Y\to X\times Y$ 使得 $(f\times g)(x,y)=(f(x),g(y))$.
        \end{enumerate}
    \end{proposition}
    \begin{proof}
        (1)\ (2) 请自行证明， (2) 的提示是：覆盖的并是并的覆盖.

        (3) 在 $m>0$ 时，$d_n^{f^m}\leq d_{nm-m+1}^f$，则 $B_f(x,\varepsilon,mn-m+1)\subset B_{f^m}(x,\varepsilon,n)$，于是 $S_d(f^m,\varepsilon,n)\leq S_d(f,\varepsilon,mn)$，于是 $\leq$ 方向得证. 另一方面，由于 $f,\ldots,f^{m-1}$ 都是一致连续的，于是 $\forall\; \varepsilon>0$ 都存在 $\delta>0$ 使得 $B(x,\delta)\subset B_f(x,\varepsilon,m)$ 对任意 $x\in X$ 成立，于是根据定义，$B_{f^m}(x,\delta,n)\subset B_f(x,\varepsilon,mn)$，类似地可以证明 $\geq$ 方向.

        (3) 在 $m<0$ 时，只需验证 $m=-1$ 的情况，这时由于 $B_f(x,\varepsilon,n)=B_{f^{-1}}(f^{n-1}(x),\varepsilon,n)$ 即可以类似证明.

        对于 (4)，首先给出乘积度量 $d((x_1,y_1),(x_2,y_2))=\max(d_X(x_1,x_2),d_Y(y_1,y_2))$（请验证其是度量且与乘积拓扑相合），于是乘积度量的开球就是开球的乘积，于是分别使用 $S_d,N_d$ 可以分别证明 $\leq,\geq$ 方向.
    \end{proof}

    \begin{theorem}
        设流 $\Phi={\{\varphi^t\}}_{t\in\R}$，则 $h(\Phi)=h(\varphi^1)$.
    \end{theorem}
    \begin{proof}
        由一致连续性我们知道 $\forall\;\varepsilon>0,\; \exists\;\delta>0,\; d(x,y)\leq\delta\Rightarrow d_1^{\Phi}(x,y)<\varepsilon$，则 $B_{\Phi}(x,\varepsilon,T)\supset B_{\varphi^1}(x,\delta,[T])$（$[T]$ 表示不大于 $T$ 的最大整数），则 $h(\Phi)\leq h(\varphi^1)$. 另一方面，由度量的定义知 $d_n^{\Phi}\geq d_n^{\varphi^1}$，则 $\geq$ 方向即证.
    \end{proof}
    \begin{remark}
        这说明流的拓扑熵可以转化为离散迭代的拓扑熵.
    \end{remark}

    在有良好的体形式定义的底空间（光滑空间）中，我们也可以考虑像的体积增长作为一种拓扑不变量（特别的，$2$ 维流形上的 $C^{\infty}$ 微分同胚通过体积增长定义的不变量与拓扑熵相等）.

    我们还可以通过代数拓扑理论衡量基本群/同调群的增长，并定义基本群熵（不大于拓扑熵）和同调熵.

    最后，我们要说明拓扑熵是一个足够好的不变量指标，也就是说，它不会轻易膨胀到无穷.

    \begin{definition}[球维数]
        设 $(X,d)$ 是紧度量空间，$b(\varepsilon)$ 是 $X$ 使用半径为 $\varepsilon$ 的球进行的最小覆盖数，则定义球维数 $D(X)=\varlimsup\limits_{\varepsilon\to 0}\frac{\log b(\varepsilon)}{\abs{\log\varepsilon}}\in\R_{\geq 0}\cup \{\infty\}$.
    \end{definition}
    \begin{remark}
        \begin{itemize}
            \item $D(X)$ 在双 Lipschitz 等价变换下不变；
            \item $D(X)$ 在微分流形上球维数就是拓扑维数；
            \item $Y\subset X\Rightarrow D(Y)\leq D(X)$；
            \item $D(X)$ 不是拓扑不变量.
        \end{itemize}
    \end{remark}

    \begin{theorem}[拓扑熵有限性]
        设 $(X,d)$ 是有有限球维数 $D(X)$ 的紧度量空间，$f:X\to X$ 是 Lipschitz 连续映射，则 $h(f)\leq D(X)\max(0,\log L(f))$，其中 $L(f)=\sup\limits_{x\neq y}\frac{d(f(x),f(y))}{d(x,y)}$ 是 Lipschitz 常数.
    \end{theorem}
    \begin{proof}
        任取 $L>\max(1,L(f))$，则 $d(f(x),f(y))\leq Ld(x,y)$，则 $f^m(B_d(x,\varepsilon))\subset B_d(f^m(x),L^m \varepsilon)$，则 $B_d(x,L^{-n}\varepsilon)\subset B_{f}(x,\varepsilon,n)$，于是 $S(f,\varepsilon,n)\leq b(L^{-n}\varepsilon)$，则 $h(f,\varepsilon)=\varlimsup\limits_{n\to\infty}\frac{1}{n}\log S(f,\varepsilon,n)\leq \varlimsup\limits_{n\to\infty}\frac{\abs{\log(L^{-n}\varepsilon)}}{n}\frac{\log b(L^{-n}\varepsilon)}{\abs{\log(L^{-n}\varepsilon)}}=D(X)\log L$，两边令 $\varepsilon\to 0$ 即证.
    \end{proof}
    \begin{remark}
        这一定理给出了映射拓扑熵的估计，这里对于 $C^1$ 映射，$L(f)=\sup\limits_{x\in X}\norm{Df_x}$，其中 $\norm{\cdot}$ 取得是谱半径.
    \end{remark}

    \section{拓扑熵的计算}
    \begin{example}[环面平移]
        平移 $T_{\gamma}$ 的拓扑熵是 $0$（对应的，线性流也是），因为 $d_n^{T_{\gamma}}=d$，使用分离集即知.
    \end{example}

    \begin{example}[均匀扩张映射]
        $h(E_m)=\log\abs{m}=p(E_m)$，这里 $p(E_m)$ 是周期点渐近增长指数. 后者容易计算出来，前者我们先分析度量（不妨设 $m>0$，因为 $m<0$ 同理）：只要 $d(x,y)<\frac{m^{-n}}{2}$，就有 $d_n^{E_m}(x,y)=d(E_m^{n-1}(x),E_m^{n-1}(y))=m^{n-1}d(x,y)$. 则在 $d(x,y)<\frac{m^{-n}}{2}$ 的前提下，$d_n^{E_m}(x,y)\geq \varepsilon$（$\varepsilon$ 充分小）当且仅当 $d(x,y)>\varepsilon m^{-n+1}$. 取 $\varepsilon=m^{-k}$，则 $\{im^{-n+1-k}\mid i\in\Z\}$ 是最大的 $(n,\varepsilon)$-分离集，于是 $N_d(E_m,m^{-k},n)=m^{n+k-1}$，则 $h(E_m)=\log m$.
    \end{example}
    \begin{remark}
        由于同度数 $m$ 的扩张映射之间拓扑共轭，因此它们的拓扑熵都是 $\log\abs{m}$.
    \end{remark}

    \begin{example}[拓扑 Markov 链]
        我们考虑 $\Omega_N$ 空间赋予的 $10N$ 度量，则对称圆柱 $C_{\alpha}^m$ 是半径为 $\varepsilon_m=\frac{{(10N)}^{-m}}{2}$ 的闭球. 类似的，$C_{\alpha}^{-m,\ldots,n+m}$ 是 $d_n^{\sigma_N}$ 的半径为 $\varepsilon_m$ 的闭球. 考虑矩阵 $A$，则我们只需数 $C_{\alpha}^{-m,\ldots,n+m}\cap \Omega_A$ 中不同 $\alpha$ 的元素，而连续秩为 $n$ 的圆柱，不同 $\alpha$ 的元素数量就是 $\sum\limits_{i,j=0}^{N-1}a_{ij}^n=\norm{A^n}_1$，也即 $N_{d_{10N}}(\sigma_A,\varepsilon_m,n)=\sum\limits_{i,j=0}^{N-1}a_{ij}^{n+2m}=\norm{A^{n+2m}}_1$，于是 $h(\sigma_A)=\lim\limits_{m\to\infty}\lim\limits_{n\to\infty}\frac{1}{n}\log\norm{A^{n+2m}}_1=\lim\limits_{m\to\infty}\lim\limits_{n\to\infty}\frac{n+2m}{n}\log\norm{A^{n+2m}}_1^{\frac{1}{n+2m}}=\lim\limits_{m\to\infty}1\cdot \log\abs{\lambda_A^{\max}}=\log\abs{\lambda_A^{\max}}$，其中 $\abs{\lambda_A^{\max}}$ 是谱半径，也就是模长最长的特征值模长.
    \end{example}
    \begin{remark}
        取 $A$ 为全 $1$ 矩阵则知 $h(\sigma_N)=\log N$.
    \end{remark}

    \begin{example}[双曲环面自同构]
        对于 $L=\begin{bmatrix}
            2&1\\
            1&1\\
        \end{bmatrix}$ 在 $\T^2$ 上诱导的双曲环面自同构，我们知道 $L$ 是 $\sigma_A$ 的因子，其中 $A=\begin{bmatrix}
            1&1&0&1&0\\
            1&1&0&1&0\\
            1&1&0&1&0\\
            0&0&1&0&1\\
            0&0&1&0&1\\
        \end{bmatrix}$，于是 $h(L)\leq h(\sigma_A)=\log\frac{3+\sqrt{5}}{2}$.

        接下来考虑周期点，首先由于特征值 $\lambda_1=\frac{3+\sqrt{5}}{2}\in (2,3)$，于是我们考虑 $\varepsilon=\frac{1}{6}$ 下的 $N_d(L,\varepsilon,n)$，其中 $n$ 充分大. 任取两个周期为 $n$ 的周期点 $p,q$，若 $d(p,q)\geq \varepsilon$ 则已互相分离，而 $d(p,q)<\varepsilon$ 时，由于 $\varepsilon$ 足够小，我们可以找到一个以特征方向为边方向的最小矩形 $psqt$（逆时针），而该矩形作用 $n$ 次后的半径会远大于该矩形本身的半径，因此不再是最小矩形，虽然 $p,q$ 两点不变，因此这时在 $\R^2$ 上 $d(L^n(p),L^n(s))>1$，从而由于在 $\R^2$ 上 $d(L^n(p),L^n(s))=\lambda_1^n d(p,s)$，因此存在 $k\in [0,n-1]$ 使得 $d(L^k(p),L^k(s))\in [\varepsilon,3\varepsilon]$，于是 $d(L^k(p),L^k(q))\in [\varepsilon,3\sqrt{2}\varepsilon]$ 在 $\R^2$ 上，由于 $3\sqrt{2}\varepsilon=\frac{1}{\sqrt{2}}$ 则这一距离也是 $\T^2$ 上的距离. 综上，$N_d(L,\varepsilon,n)\geq P_n(L)=\lambda_1^n+\lambda_1^{-n}-2$，于是 $h(L)\geq h_d(L,\varepsilon)=\varliminf\limits_{n\to\infty}\frac{1}{n}\log N_d(L,\varepsilon,n)=\log \lambda_1=\log \frac{3+\sqrt{5}}{2}$. 于是 $h(L)=\log\frac{3+\sqrt{5}}{2}$.
    \end{example}

\end{document}